\documentclass[12pt]{article}

\usepackage[a4paper, margin=2.5cm]{geometry}
\usepackage{amsmath}
\usepackage{amssymb}
\usepackage{amsthm}
\usepackage[colorlinks=true, linkcolor=blue, citecolor=blue, urlcolor=blue]{hyperref}
\usepackage{booktabs}
\usepackage{natbib}

\newtheorem{theorem}{Theorem}
\newtheorem{proposition}[theorem]{Proposition}
\newtheorem{definition}[theorem]{Definition}
\newtheorem{remark}[theorem]{Remark}

\newcommand{\Mcal}{\mathcal{M}}
\newcommand{\Acal}{\mathcal{A}}
\newcommand{\phig}{\varphi}

\title{From Class Energy to Mass:\\The Nonlinear Map Problem for Closure-Class Particles}
\author{%
  Lee Smart\\[4pt]
  \textit{Institute of Vibrational Field Dynamics}\\[2pt]
  \texttt{contact@vibrationalfielddynamics.org}\\[2pt]
  \texttt{@vfd\_org}%
}
\date{April 2026}

\begin{document}

\maketitle

\noindent\textbf{Status:} Preprint (not peer-reviewed)

\begin{abstract}
Companion papers established that particles correspond to stable closure classes on a $\phig$-structured manifold (Paper~II) and that the manifold metric induces class energies with ratios of order $10^1$--$10^2$ (Paper~III). The observed proton-to-electron mass ratio of $\sim\!1836$ requires a strongly nonlinear map from class energy to physical mass. This paper evaluates seven candidate map families against the established closure-class structure. Two zero-parameter maps produce phi-integer-power ratios that bracket the observed value: the combinatorial map $\phig^{\Delta C} = \phig^{15} \approx 1364$ (derived from shell complexity) and the multiplicative shell product $\phig^{16} \approx 2207$ (derived from metric eigenvalue products). The observed ratio at $\phig^{15.618}$ falls between these bounds, identifying the inter-shell composition rule --- the precise manner in which occupied shells combine to determine mass --- as the primary remaining open problem. The paper is presented as a structural diagnosis, not a mass derivation.
\end{abstract}

%==================================================================
\section{Introduction}
%==================================================================

The preceding papers in this sequence established:
\begin{itemize}
    \item Paper~I \citep{Smart2026}: the crystallisation operator, a deterministic mechanism for state selection.
    \item Paper~II: closure classes as particle identities, with discrete labels $(S, w, \sigma)$ and combinatorial invariant $\Delta C = 15$ between proton and electron.
    \item Paper~III: five candidate metric families on the $\phig$-manifold, with metric-induced class-energy ratios of order $10^1$--$10^2$.
\end{itemize}

Paper~III diagnosed that the additive class-energy ratio ($\sim\!27$ under the minimal metric) falls far short of the observed mass ratio ($\sim\!1836$). The relationship between geometric class energy and physical mass must therefore be nonlinear.

This paper addresses the resulting question: \textit{what nonlinear map transforms class energy into mass?}

\paragraph{Scope.} Seven candidate map families are evaluated. No map is fitted to the observed ratio. The paper identifies the viable family and the remaining open structure.

%==================================================================
\section{Why a Nonlinear Map Is Required}\label{sec:why}
%==================================================================

Under the minimal radial metric (Paper~III, MF1), the spectral class energies are:
\begin{align}
    E_{\mathrm{class}}(\text{electron}) &= \phig^2 \approx 2.618 \\
    E_{\mathrm{class}}(\text{proton}) &= \phig^4 + \phig^6 + \phig^8 \approx 71.78
\end{align}

The additive ratio is $E_p / E_e \approx 27.4$. The observed mass ratio is $m_p / m_e \approx 1836.15$. The discrepancy is a factor of $\sim\!67$. No choice of metric closes this gap while remaining additive; the map must involve exponential or multiplicative amplification.

%==================================================================
\section{Candidate Map Families}\label{sec:maps}
%==================================================================

Seven map families are evaluated. Each maps closure-class data $(S, E_{\mathrm{class}}, C)$ to a mass surrogate. The ratio $m_p/m_e$ cancels any universal calibration constant.

\subsection{F2: Combinatorial Map ($\phig^C$)}

\begin{equation}
    m(\Acal) \sim \phig^{C(\Acal)}
\end{equation}

Mass scales exponentially with the combinatorial complexity $C = \prod n_i - \sum n_i - 1$. This is Paper~II's original hypothesis. It uses only the shell support, not the metric energy.

Ratio: $\phig^{15} \approx 1364$. Error: $25.7\%$.

\textit{Origin:} Postulated in Paper~II. Structurally motivated but not derived from the metric.

\subsection{F3: Multiplicative Shell Product}

\begin{equation}\label{eq:F3}
    m(\Acal) \sim \prod_{n \in S} \lambda(n) = \prod_{n \in S} \phig^{2n}
\end{equation}

Mass is the \textit{product} of per-shell spectral weights, not their sum. Each shell contributes multiplicatively to the total bound-state mass.

Ratio: $\phig^{18}/\phig^2 = \phig^{16} \approx 2207$. Error: $20.2\%$.

\textit{Origin:} Derived. The spectral weights $\lambda(n) = \phig^{2n}$ come from the manifold metric (Paper~III, MF1). The product is the natural multiplicative composition.

\subsection{F5: Mode-Weighted Phi Product}

\begin{equation}
    m(\Acal) \sim \prod_{n \in S} n \cdot \phig^n
\end{equation}

Each shell contributes its index (mode count) times its $\phig$-weight.

Ratio: $\approx 1127$. Error: $38.6\%$.

\textit{Origin:} Derived. Combines mode multiplicity with geometry.

\subsection{Other Families Tested}

\begin{itemize}
    \item \textbf{F1} (pure exponential $e^{E/E_0}$): overflow --- exponential of energy is too strong.
    \item \textbf{F4} ($\phig^{E_{\mathrm{class}}}$): overflow --- energy-as-exponent diverges for proton.
    \item \textbf{F6} (quadratic closure product $\prod \phig^{n^2}$): $\phig^{28} \approx 7 \times 10^5$ --- far too large.
    \item \textbf{F7} ($\phig^C \times E_{\mathrm{class}}$): $\sim\!37400$ --- overshoots by factor $\sim\!20$.
\end{itemize}

All four are rejected on structural grounds (overflow, wrong scale, or combination of independently valid quantities yielding implausible results).

%==================================================================
\section{Results}\label{sec:results}
%==================================================================

\begin{table}[ht]
\centering
\caption{Viable mass-map families (zero free parameters).}
\label{tab:results}
\begin{tabular}{@{}llccl@{}}
\toprule
Family & Map & Predicted $m_p/m_e$ & Error & Origin \\
\midrule
F2 & $\phig^{\Delta C} = \phig^{15}$ & 1364 & 25.7\% & Postulated (combinatorial) \\
F3 & $\prod \lambda(n) \to \phig^{16}$ & 2207 & 20.2\% & Derived (metric eigenvalue product) \\
F5 & $\prod n \cdot \phig^n$ & 1127 & 38.6\% & Derived (mode $\times$ geometry) \\
\bottomrule
\end{tabular}
\end{table}

The observed ratio $m_p/m_e = 1836.15$ corresponds to $\phig^{15.618}$. It falls between the two strongest zero-parameter predictions:
\begin{equation}
    \phig^{15} \;<\; m_p/m_e \;<\; \phig^{16}
\end{equation}

Neither the combinatorial map (F2, exponent $15$) nor the multiplicative metric map (F3, exponent $16$) is exact. The observed exponent $15.618$ is not an integer, nor a simple $\phig$-expression.

%==================================================================
\section{The Bracketing Result}\label{sec:bracketing}
%==================================================================

\begin{proposition}[Bracketing]\label{prop:bracket}
Under the established closure-class assignments and the minimal $\phig$-metric:
\begin{enumerate}
    \item The combinatorial map $m \sim \phig^C$ produces $m_p/m_e = \phig^{15}$ (lower bound).
    \item The multiplicative spectral map $m \sim \prod \lambda(n)$ produces $m_p/m_e = \phig^{16}$ (upper bound).
    \item Both are derived from the framework with zero free parameters.
    \item The observed ratio $\phig^{15.618}$ lies strictly between them.
\end{enumerate}
\end{proposition}

The combinatorial map uses only shell-index arithmetic; the multiplicative map uses the metric eigenvalues. Their difference (exponent $15$ vs $16$) reflects a single unit of ``metric excess'' --- the multiplicative map assigns more weight to higher shells than the combinatorial map does.

The physical mass map must interpolate between these bounds. The interpolation rule is determined by how occupied shells compose --- whether fully multiplicatively (F3), purely combinatorially (F2), or through a partial composition that reflects inter-shell coupling.

%==================================================================
\section{The Inter-Shell Composition Rule}\label{sec:composition}
%==================================================================

The missing object is the \textit{inter-shell composition rule}: the exact prescription for how the spectral weights of occupied shells combine to determine mass.

\begin{itemize}
    \item \textbf{Additive composition} ($E = \sum \lambda_n$): gives ratio $\sim\!27$. Too weak.
    \item \textbf{Fully multiplicative composition} ($m = \prod \lambda_n$): gives $\phig^{16}$. Too strong by $20\%$.
    \item \textbf{Combinatorial composition} ($m = \phig^C$): gives $\phig^{15}$. Too weak by $26\%$.
\end{itemize}

The observed ratio lies between the multiplicative and combinatorial extremes. This suggests a \textit{partial multiplicative} composition in which not all shells contribute with their full spectral weight.

Possible mechanisms for partial composition:
\begin{enumerate}
    \item Adjacent shells interact multiplicatively, but non-adjacent shells do not (partial coupling).
    \item Shell contributions are weighted by a function intermediate between $\phig^{2n}$ (full metric) and the combinatorial index (pure arithmetic).
    \item The boundary shell (electron, $n=1$) has a distinct composition rule from interior shells.
\end{enumerate}

Determining this rule requires either a derivation from the closure dynamics or the manifold geometry, or additional physical input. This is the primary open problem.

%==================================================================
\section{Relation to $\Delta C$}\label{sec:deltaC}
%==================================================================

Paper~II's combinatorial invariant $\Delta C = 15$ remains structurally significant:
\begin{itemize}
    \item It provides the lower bound $\phig^{15}$ on the mass ratio.
    \item The multiplicative map exponent ($16 = 2 \times \sum_p n_i - 2 \times \sum_e n_i = 2 \times 9 - 2 \times 1$) is a distinct quantity.
    \item The observed exponent $15.618$ is closer to $15$ than to $16$, suggesting the combinatorial structure dominates with a metric correction of $+0.618 \approx 1/\phig$.
\end{itemize}

Whether the correction $+1/\phig$ has structural origin or is coincidental cannot be determined from the map analysis alone.

%==================================================================
\section{Implications for Neutron-Proton Splitting}\label{sec:neutron}
%==================================================================

All nonlinear map families amplify small differences in class energy. If the neutron's class energy differs from the proton's by a symmetry-sector or torsional correction $\delta E$, the mass splitting is:
\begin{equation}
    \Delta m \sim m_p \cdot \frac{\delta E}{E_p} \cdot \alpha_{\mathrm{map}}
\end{equation}
where $\alpha_{\mathrm{map}}$ is the map's amplification factor at the proton class energy. For the multiplicative map, the amplification is logarithmic in the shell weights; for the combinatorial map, it depends on the sensitivity of $C$ to perturbation.

The neutron-proton mass difference ($\Delta m / m_p \approx 0.14\%$) is plausible under any of the viable map families, provided the symmetry-sector correction $\delta E$ is of order $0.1$--$1\%$ of the proton class energy.

%==================================================================
\section{Discussion}\label{sec:discussion}
%==================================================================

\subsection{What This Paper Establishes}

\begin{itemize}
    \item A nonlinear class-energy-to-mass map is required.
    \item Two zero-parameter maps bracket the observed ratio: $\phig^{15} < m_p/m_e < \phig^{16}$.
    \item The multiplicative shell product (F3) is the most structurally justified nonlinear map.
    \item The inter-shell composition rule is the primary remaining open problem.
\end{itemize}

\subsection{What This Paper Does Not Establish}

\begin{itemize}
    \item The exact mass ratio. The best zero-parameter prediction is $\phig^{16}$ (20\% error).
    \item The inter-shell composition rule.
    \item The neutron-proton mass difference (only assessed as plausible, not derived).
\end{itemize}

\subsection{The Four-Layer Architecture}

\begin{table}[ht]
\centering
\begin{tabular}{@{}lll@{}}
\toprule
Layer & Paper & Status \\
\midrule
Selection dynamics & Paper I & Established \\
Class structure & Paper II & Established \\
Manifold geometry & Paper III & Established \\
Nonlinear mass map & Paper IV (this paper) & Bracketed; composition rule open \\
\bottomrule
\end{tabular}
\end{table}

%==================================================================
\section{Conclusion}
%==================================================================

The class-energy-to-mass map must be nonlinear. Two structurally justified maps with zero free parameters bracket the observed proton-to-electron mass ratio: the combinatorial map $\phig^{15}$ from below and the multiplicative spectral product $\phig^{16}$ from above. The observed ratio at $\phig^{15.618}$ falls between these bounds, identifying the inter-shell composition rule as the final missing mathematical object. The four-layer architecture --- selection, classification, geometry, and mass map --- is now established through three complete layers and one partially constrained layer. The remaining open problem is precise and well-defined: how do occupied shells compose to determine mass?

%==================================================================
\bibliographystyle{plainnat}
\bibliography{references}

\end{document}
